\section{الفصل~\ref{sec:motorlearning}. تعلّم الحركة}

قاعدة المربعات الصغرى وفائدة سلك الخطأ المستقل مبرهنتان؛ وبنية الدارة ذات سلك الخطأ، والإضعاف عند التزامن، وضبط الأثر على التأخير، والأثر اللاحق، والعودة التلقائية للمهارة مقيسة؛ أما أن الدارة كلها تعمل تعلمًا موجَّهًا وتبني نموذجًا داخليًا ففرضية رائدة؛ وأرقام نموذج العمليتين حساب بنموذج الكتاب.

{\footnotesize
\setlength{\tabcolsep}{3pt}\renewcommand{\arraystretch}{1.15}
\begin{longtable}{>{\raggedright}p{0.5cm}>{\raggedright}p{4.0cm}>{\raggedright}p{5.6cm}>{\raggedright}p{2.3cm}>{\raggedright\arraybackslash}p{1.7cm}}
\toprule
م & القضية & التجربة وما قيس & المصدر & الحالة \\
\midrule\endfirsthead
\toprule
م & القضية & التجربة وما قيس & المصدر & الحالة \\
\midrule\endhead
1 & القاعدة~\eqref{eq:lms} تسير في المتوسط مع تدرج مربع الخطأ & نتيجة من نظرية المرشحات التكيّفية: التصحيح المتناسب مع الدخل وخطأ المخرج انحدار تدرجي عشوائي لمتوسط مربع الخطأ. & \cite{WidrowHoff1960} & نظرية \\
2 & مع سلك خطأ لكل مخرج لا تتوقف السرعة على عدد المخارج؛ ومع عدد مشترك واحد تنخفض مع عدد العصبونات الضاجّة (القضية~\ref{prop:teachers}) & منحنيات التعلم لشبكة خطية: التعلم باضطرابات عشوائية للعصبونات أبطأ من التدرج الدقيق بعدد مرات يساوي عدد العصبونات المضطربة. & \cite{Werfel2005} & نظرية \\
3 & الدارة ذات سلك الخطأ تعمل تعلمًا موجَّهًا (القضية~\ref{prop:errorwire}) & نظريات مستنتجة من بنية الدارة. تجارب الأسطر من 7 إلى 10 تتفق معها؛ أما أن الدارة كلها تعمل هكذا بالذات فلم يُبرهن. & \cite{Marr1969, Albus1971} & فرضية \\
4 & طبقة التوسيع: نحو $7\cdot10^{10}$ عصبون \nt{GLUT} صغير، أكثر مما في بقية الدماغ كله & عدّ النوى في معلّق من النسيج المذاب: في مخيخ الإنسان $69\cdot10^9$ عصبونًا من أصل $86\cdot10^9$ في الدماغ كله. القياس المجسّم: $101\cdot10^9$ عصبون صغير في مخيخ رجال مسنين. & \cite{Azevedo2009, Andersen1992cereb} & مُقاس \\
5 & رفع بُعد الدخل يجعل العلاقة المطلوبة خطية & مبرهنة عدد التقسيمات: المجموع الموزون لـ$n$ مدخلًا مع عتبة ينفذ أي تصنيف لنحو $2n$ متجهًا في وضع عام. أما أن المخيخ يستعمل هذا بالذات فجزء من فرضية السطر~3. & \cite{Cover1965} & نظرية \\
6 & نحو $3\cdot10^7$ عصبون \nt{GABA} مخرج، لكل منها نحو $10^5$ مدخل & القياس المجسّم لمخيخ الإنسان: $30.5\cdot10^6$ عصبون مخرج. المجهر الإلكتروني، الجرذ: نحو $1.75\cdot10^5$ مدخل من طبقة التوسيع لكل عصبون مخرج. والناقل العصبي المثبِّط لعصبونات المخرج في المراجعة. & \cite{Andersen1992cereb, NapperHarvey1988, Ito2001} & مُقاس \\
7 & سلك خطأ واحد بالضبط لكل عصبون مخرج، نحو مرة في الثانية، وفي كل مرة ومضة قوية & مراجعة للتسجيلات: كل عصبون مخرج يتلقى مدخل \nt{GLUT} متسلقًا واحدًا، يستدعي نبضة معقدة بتردد من رتبة $1\,\text{Hz}$. & \cite{Ito2001} & مُقاس \\
8 & احتمال الومضة يتوقف على اتجاه خطأ الحركة & القرد، الجزء المحرك للعين في المخيخ، تكيّف الرمشات السريعة: اتجاه الخطأ البصري يحدد احتمال النبضة المعقدة، وحجمه توقيتها؛ والومضة تغيّر النبضات البسيطة في المحاولة التالية وتزيح الحركة على امتداد الخطأ المفضل للخلية. & \cite{Herzfeld2018} & مُقاس \\
9 & تزامن ومضة الخطأ مع مدخل نشط يُضعف المدخل ($-\eta e_i x_j$) & مراجعة للتجارب: التنبيه المشترك لمدخل من طبقة التوسيع ولسلك الخطأ يُحدث إضعافًا طويل الأمد لهذا المدخل؛ وتنبيه المدخل وحده لا يُحدثه. & \cite{Ito2001} & مُقاس \\
10 & طول الأثر مضبوط على تأخير الخطأ: عند تأخير $100\ms$ يضعف أكثر ما يضعف المدخل الذي كان نشطًا قبله بـ$100\ms$ & شرائح وفئران حية: في الجزء المسؤول عن تعلم حركة العين يكون الإضعاف مضبوطًا بدقة على فاصل نحو $120\ms$ بين المدخل والخطأ، وهو زمن وصول إشارة الخطأ في هذه المهمة؛ وفي جزء آخر تُضبط خلايا مختلفة على فواصل مختلفة. & \cite{Suvrathan2016} & مُقاس \\
11 & الدارة مرشح تكيّفي~\eqref{eq:adaptivefilter} يتعلم نموذجًا داخليًا للجسم & نموذج مستنتج من بنية الدارة. ولا توجد تجربة مباشرة قيس فيها المرشح المتعلَّم بوصفه دالة تحويل الجسم. & \cite{Fujita1982} & فرضية \\
12 & التنبؤ يطرح نتائج الحركات الذاتية، ولهذا لا ندغدغ أنفسنا & بشر: اللمسة الذاتية تبدو أقل دغدغة من اللمسة نفسها من الغير؛ وفي التصوير تستجيب القشرة الحسية الجسدية للمسة الذاتية أضعف، ويقل نشاط المخيخ حين تلمس الحركة الجلد. والتفسير بطرح التنبؤ فرضية. & \cite{Blakemore1998} & فرضية \\
13 & مع قوة خارجية يتناقص الخطأ، وبعد إزالتها تخطئ اليد في الاتجاه الآخر & بشر يمدّون أيديهم بمقبض روبوت في حقل قوة: تعود المسارات مستقيمة؛ وبعد إزالة الحقل يكون الأثر اللاحق شبه صورة مرآوية للتشوهات الأولى؛ وينتقل أيضًا إلى أجزاء من الفضاء لم يُتدرَّب فيها. & \cite{ShadmehrMussaIvaldi1994} & مُقاس \\
14 & تتعلم عمليتان~\eqref{eq:tworate}؛ وبعد الاضطراب المعاكس تعود المهارة بنفسها & بشر، حقل قوة، ثم حقل معاكس قصير ومحاولات مع تثبيت الخطأ: يعود التصحيح بنفسه إلى القديم؛ ونموذج العمليتين يعيد إنتاج ذلك، ونموذج العملية الواحدة لا. & \cite{Smith2006} & مُقاس \\
15 & أرقام الشكل~\ref{fig:tworate}: $0.84$ (البطيئة $0.63$) خلال $200$ حركة؛ $6$ معاكسة؛ السريعة $-0.53$ والبطيئة $0.46$؛ عودة إلى $0.37$ خلال $40$ حركة & حساب بـ\foreignlanguage{english}{\texttt{models/\allowbreak motor\_\allowbreak learning.py}}، $A_f=0.92$، $B_f=0.1$، $A_s=0.996$، $B_s=0.02$: $0.840$ ($0.634$ و$0.205$)، $6$ حركات، $-0.531$ و$0.457$، وقمة $0.371$ عند الحركة الـ$40$. & \foreignlanguage{english}{\texttt{models/\allowbreak motor\_\allowbreak learning.py}} & نموذج \\
16 & تلزم عمليتان لأن الجسم يتغير بسرعات مختلفة & نظرية: التقدير الأمثل مع اضطرابات بأعمار متعددة يعطي عدة مرشحات بثوابت زمنية مختلفة، ويفسر العودة التلقائية وتسارع إعادة التعلم. & \cite{Kording2007} & فرضية \\
\bottomrule
\end{longtable}}
